\documentclass[11pt]{article}

% ---- theme variables (passed via pandoc -V) ----
\newcommand{\ThemeName}{System Design Problems — Advanced}
\newcommand{\ThemeKicker}{Design Track}
\newcommand{\ThemeNumber}{II}

\usepackage{interviewstyle}
\setthemecolor{db2777}{f472b6}

% pandoc-required plumbing
\providecommand{\tightlist}{\setlength{\itemsep}{1pt}\setlength{\parskip}{0pt}}
\setlength{\emergencystretch}{3em}
\providecommand{\pandocbounded}[1]{#1}

% syntax-highlighting macros emitted by pandoc (skylighting)
\usepackage{color}
\usepackage{fancyvrb}
\newcommand{\VerbBar}{|}
\newcommand{\VERB}{\Verb[commandchars=\\\{\}]}
\DefineVerbatimEnvironment{Highlighting}{Verbatim}{commandchars=\\\{\}}
% Add ',fontsize=\small' for more characters per line
\usepackage{framed}
\definecolor{shadecolor}{RGB}{35,38,41}
\newenvironment{Shaded}{\begin{snugshade}}{\end{snugshade}}
\newcommand{\AlertTok}[1]{\textcolor[rgb]{0.58,0.85,0.30}{\textbf{\colorbox[rgb]{0.30,0.12,0.14}{#1}}}}
\newcommand{\AnnotationTok}[1]{\textcolor[rgb]{0.25,0.50,0.35}{#1}}
\newcommand{\AttributeTok}[1]{\textcolor[rgb]{0.16,0.50,0.73}{#1}}
\newcommand{\BaseNTok}[1]{\textcolor[rgb]{0.96,0.45,0.00}{#1}}
\newcommand{\BuiltInTok}[1]{\textcolor[rgb]{0.50,0.55,0.55}{#1}}
\newcommand{\CharTok}[1]{\textcolor[rgb]{0.24,0.68,0.91}{#1}}
\newcommand{\CommentTok}[1]{\textcolor[rgb]{0.48,0.49,0.49}{#1}}
\newcommand{\CommentVarTok}[1]{\textcolor[rgb]{0.50,0.55,0.55}{#1}}
\newcommand{\ConstantTok}[1]{\textcolor[rgb]{0.15,0.68,0.68}{\textbf{#1}}}
\newcommand{\ControlFlowTok}[1]{\textcolor[rgb]{0.99,0.74,0.29}{\textbf{#1}}}
\newcommand{\DataTypeTok}[1]{\textcolor[rgb]{0.16,0.50,0.73}{#1}}
\newcommand{\DecValTok}[1]{\textcolor[rgb]{0.96,0.45,0.00}{#1}}
\newcommand{\DocumentationTok}[1]{\textcolor[rgb]{0.64,0.20,0.25}{#1}}
\newcommand{\ErrorTok}[1]{\textcolor[rgb]{0.85,0.27,0.33}{\underline{#1}}}
\newcommand{\ExtensionTok}[1]{\textcolor[rgb]{0.00,0.60,1.00}{\textbf{#1}}}
\newcommand{\FloatTok}[1]{\textcolor[rgb]{0.96,0.45,0.00}{#1}}
\newcommand{\FunctionTok}[1]{\textcolor[rgb]{0.56,0.27,0.68}{#1}}
\newcommand{\ImportTok}[1]{\textcolor[rgb]{0.15,0.68,0.38}{#1}}
\newcommand{\InformationTok}[1]{\textcolor[rgb]{0.77,0.36,0.00}{#1}}
\newcommand{\KeywordTok}[1]{\textcolor[rgb]{0.81,0.81,0.76}{\textbf{#1}}}
\newcommand{\NormalTok}[1]{\textcolor[rgb]{0.81,0.81,0.76}{#1}}
\newcommand{\OperatorTok}[1]{\textcolor[rgb]{0.81,0.81,0.76}{#1}}
\newcommand{\OtherTok}[1]{\textcolor[rgb]{0.15,0.68,0.38}{#1}}
\newcommand{\PreprocessorTok}[1]{\textcolor[rgb]{0.15,0.68,0.38}{#1}}
\newcommand{\RegionMarkerTok}[1]{\textcolor[rgb]{0.16,0.50,0.73}{\colorbox[rgb]{0.08,0.19,0.26}{#1}}}
\newcommand{\SpecialCharTok}[1]{\textcolor[rgb]{0.24,0.68,0.91}{#1}}
\newcommand{\SpecialStringTok}[1]{\textcolor[rgb]{0.85,0.27,0.33}{#1}}
\newcommand{\StringTok}[1]{\textcolor[rgb]{0.96,0.31,0.31}{#1}}
\newcommand{\VariableTok}[1]{\textcolor[rgb]{0.15,0.68,0.68}{#1}}
\newcommand{\VerbatimStringTok}[1]{\textcolor[rgb]{0.85,0.27,0.33}{#1}}
\newcommand{\WarningTok}[1]{\textcolor[rgb]{0.85,0.27,0.33}{#1}}

\begin{document}

\makecoverpage

\thispagestyle{empty}
{\hypersetup{hidelinks}\tableofcontents}
\clearpage

Four fully worked FAANG-style system-design solutions. Each covers requirements,
capacity math, API, data model, architecture, deep dives, scaling, and trade-offs.

\section{1. Design a Video Streaming Service (YouTube / Netflix)}\label{1-design-a-video-streaming-service-youtube--netflix}

\subsection{1. Requirements}\label{1-requirements}

\textbf{Clarifying questions to ask:}

\begin{itemize}
\tightlist
\item
  Are we designing upload + playback, or playback-only?
\item
  Live streaming, or VOD (video on demand), or both?
\item
  What\textquotesingle s the target geography --- global CDN or single-region?
\item
  Should we design the recommendation engine, or just the data hooks?
\item
  Premium / DRM requirements?
\end{itemize}

\textbf{Functional Requirements:}

\begin{itemize}
\tightlist
\item
  Users can upload videos (any resolution, up to \textasciitilde10 GB raw).
\item
  Platform transcodes video to multiple resolutions (240p, 360p, 720p, 1080p, 4K).
\item
  Users can stream videos with adaptive bitrate (quality switches based on bandwidth).
\item
  Search, browse, and play videos; like, comment, subscribe.
\item
  Recommendation feed on homepage.
\end{itemize}

\textbf{Non-Functional Requirements:}

\begin{itemize}
\tightlist
\item
  Availability: 99.99\% (\textless{} 53 min/yr downtime).
\item
  Low latency playback start: \textless{} 2 s p99.
\item
  Upload should not block playback --- async processing pipeline.
\item
  Globally distributed; content served from CDN edge nodes.
\item
  Durability: 11 nines (video assets are irreplaceable for creators).
\item
  Scale: YouTube-tier --- 500 hours of video uploaded per minute; 1B+ DAU.
\end{itemize}

\subsection{2. Capacity Estimation}\label{2-capacity-estimation}

\textbf{Users \& traffic:}

\setcodelang{Python}

\begin{Shaded}
\begin{Highlighting}[]
\NormalTok{DAU             }\OperatorTok{=} \DecValTok{1}\NormalTok{ billion}
\NormalTok{Avg watch time  }\OperatorTok{=} \DecValTok{30} \BuiltInTok{min}\OperatorTok{/}\NormalTok{day}
\NormalTok{Videos uploaded }\OperatorTok{=} \DecValTok{500}\NormalTok{ hours}\OperatorTok{/}\BuiltInTok{min} \OperatorTok{=} \DecValTok{30}\NormalTok{,}\DecValTok{000}\NormalTok{ hours}\OperatorTok{/}\BuiltInTok{min}\NormalTok{ raw footage}
\end{Highlighting}
\end{Shaded}

\textbf{Read QPS (playback requests):}

\setcodelang{}

\begin{verbatim}
1B users × 5 video views/day = 5B views/day
5B / 86,400 s ≈ 58,000 views/sec  →  ~60K QPS peak
\end{verbatim}

\textbf{Write QPS (uploads):}

\setcodelang{}

\begin{verbatim}
500 hours/min ÷ 60 = ~8.3 hours/sec raw upload
~100 concurrent upload streams sustained at any time (most uploads are large, slow)
\end{verbatim}

\textbf{Storage (per year):}

\setcodelang{}

\begin{verbatim}
Raw:         500 hr/min × 60 min/hr × 24 hr/day × 365 days = 262,800,000 hr/yr
             Avg 1 hr video raw ≈ 2 GB → 525 PB/yr raw
Transcoded:  ~5 resolutions × 0.4× compression ratio each = ~2× raw
             Total ≈ 1 EB/yr  (YouTube stores ~1 EB today)
\end{verbatim}

\textbf{Bandwidth (outbound CDN):}

\setcodelang{}

\begin{verbatim}
60K concurrent streams × avg 3 Mbps (720p) = 180 Gbps
Peak (2× avg)                               = 360 Gbps served by CDN PoPs
\end{verbatim}

\textbf{Cost hint:} At \$0.01/GB CDN egress, 180 Gbps \(\times\) 3600 s \(\times\) 24 h = \textasciitilde\$1.5 M/day --- hence CDN caching is existential.

\subsection{3. API Design}\label{3-api-design}

\setcodelang{Python}

\begin{Shaded}
\begin{Highlighting}[]
\CommentTok{\# Upload}
\NormalTok{POST   }\OperatorTok{/}\NormalTok{v1}\OperatorTok{/}\NormalTok{videos}\OperatorTok{/}\NormalTok{initiate}\OperatorTok{{-}}\NormalTok{upload}
\NormalTok{       Body: \{ title, description, tags, category \}}
\NormalTok{       Returns: \{ upload\_id, upload\_url (signed S3) \}}

\NormalTok{PUT    }\OperatorTok{/}\NormalTok{v1}\OperatorTok{/}\NormalTok{videos}\OperatorTok{/}\NormalTok{upload}\OperatorTok{/}\NormalTok{\{upload\_id\}}\OperatorTok{/}\NormalTok{chunk}
\NormalTok{       Headers: Content}\OperatorTok{{-}}\NormalTok{Range, Content}\OperatorTok{{-}}\NormalTok{Length}
\NormalTok{       Body: }\OperatorTok{\textless{}}\NormalTok{binary chunk}\OperatorTok{\textgreater{}}

\NormalTok{POST   }\OperatorTok{/}\NormalTok{v1}\OperatorTok{/}\NormalTok{videos}\OperatorTok{/}\NormalTok{upload}\OperatorTok{/}\NormalTok{\{upload\_id\}}\OperatorTok{/}\NormalTok{complete}
\NormalTok{       Returns: \{ video\_id, status: }\StringTok{"processing"}\NormalTok{ \}}

\CommentTok{\# Playback}
\NormalTok{GET    }\OperatorTok{/}\NormalTok{v1}\OperatorTok{/}\NormalTok{videos}\OperatorTok{/}\NormalTok{\{video\_id\}}
\NormalTok{       Returns: \{ metadata, manifest\_url (HLS}\OperatorTok{/}\NormalTok{DASH), thumbnail\_url \}}

\NormalTok{GET    }\OperatorTok{/}\NormalTok{v1}\OperatorTok{/}\NormalTok{videos}\OperatorTok{/}\NormalTok{\{video\_id\}}\OperatorTok{/}\NormalTok{manifest.m3u8   (served by CDN)}

\CommentTok{\# Discovery}
\NormalTok{GET    }\OperatorTok{/}\NormalTok{v1}\OperatorTok{/}\NormalTok{feed?user\_id}\OperatorTok{=\&}\NormalTok{page\_token}\OperatorTok{=}
\NormalTok{GET    }\OperatorTok{/}\NormalTok{v1}\OperatorTok{/}\NormalTok{search?q}\OperatorTok{=\&}\NormalTok{sort}\OperatorTok{=}\NormalTok{relevance}\OperatorTok{\&}\NormalTok{page}\OperatorTok{=}
\NormalTok{GET    }\OperatorTok{/}\NormalTok{v1}\OperatorTok{/}\NormalTok{videos}\OperatorTok{/}\NormalTok{\{video\_id\}}\OperatorTok{/}\NormalTok{recommendations}

\CommentTok{\# Engagement}
\NormalTok{POST   }\OperatorTok{/}\NormalTok{v1}\OperatorTok{/}\NormalTok{videos}\OperatorTok{/}\NormalTok{\{video\_id\}}\OperatorTok{/}\NormalTok{views      (fire}\OperatorTok{{-}}\KeywordTok{and}\OperatorTok{{-}}\NormalTok{forget)}
\NormalTok{POST   }\OperatorTok{/}\NormalTok{v1}\OperatorTok{/}\NormalTok{videos}\OperatorTok{/}\NormalTok{\{video\_id\}}\OperatorTok{/}\NormalTok{likes}
\NormalTok{POST   }\OperatorTok{/}\NormalTok{v1}\OperatorTok{/}\NormalTok{videos}\OperatorTok{/}\NormalTok{\{video\_id\}}\OperatorTok{/}\NormalTok{comments}
\end{Highlighting}
\end{Shaded}

\subsection{4. Data Model}\label{4-data-model}

\textbf{Metadata DB --- PostgreSQL (sharded) or Spanner (global):}

\setcodelang{SQL}

\begin{Shaded}
\begin{Highlighting}[]
\NormalTok{videos (}
\NormalTok{  video\_id        UUID }\KeywordTok{PRIMARY} \KeywordTok{KEY}\NormalTok{,}
\NormalTok{  creator\_id      BIGINT,}
\NormalTok{  title           TEXT,}
\NormalTok{  description     TEXT,}
\NormalTok{  status          ENUM(}\StringTok{\textquotesingle{}uploading\textquotesingle{}}\NormalTok{,}\StringTok{\textquotesingle{}processing\textquotesingle{}}\NormalTok{,}\StringTok{\textquotesingle{}ready\textquotesingle{}}\NormalTok{,}\StringTok{\textquotesingle{}failed\textquotesingle{}}\NormalTok{),}
\NormalTok{  duration\_s      }\DataTypeTok{INT}\NormalTok{,}
\NormalTok{  raw\_s3\_key      TEXT,}
\NormalTok{  created\_at      TIMESTAMPTZ,}
\NormalTok{  view\_count      BIGINT,     }\CommentTok{{-}{-} denormalized, updated async}
\NormalTok{  like\_count      BIGINT}
\NormalTok{)}

\NormalTok{video\_renditions (}
\NormalTok{  video\_id        UUID,}
\NormalTok{  resolution      TEXT,       }\CommentTok{{-}{-} \textquotesingle{}1080p\textquotesingle{}, \textquotesingle{}720p\textquotesingle{}, ...}
\NormalTok{  codec           TEXT,       }\CommentTok{{-}{-} \textquotesingle{}h264\textquotesingle{}, \textquotesingle{}h265\textquotesingle{}, \textquotesingle{}av1\textquotesingle{}}
\NormalTok{  manifest\_key    TEXT,       }\CommentTok{{-}{-} S3 key for HLS/DASH manifest}
\NormalTok{  segment\_prefix  TEXT,       }\CommentTok{{-}{-} S3 prefix for .ts segments}
\NormalTok{  size\_bytes      BIGINT}
\NormalTok{)}

\NormalTok{channels (channel\_id, owner\_user\_id, subscriber\_count, }\OperatorTok{..}\NormalTok{.)}
\NormalTok{subscriptions (subscriber\_id, channel\_id, created\_at)}
\NormalTok{comments (comment\_id, video\_id, user\_id, }\KeywordTok{body}\NormalTok{, created\_at)}
\end{Highlighting}
\end{Shaded}

\textbf{Why PostgreSQL/Spanner?} Transactional metadata (status updates, renditions) needs ACID. Spanner if truly global consistency is required; sharded Postgres otherwise.

\textbf{Counters --- Redis or dedicated counter service:}

\begin{itemize}
\tightlist
\item
  \texttt{VIEW:\{video\_id\}} \(\rightarrow\) HyperLogLog for approximate unique views; batched flush to Postgres.
\end{itemize}

\textbf{Search --- Elasticsearch:}

\begin{itemize}
\tightlist
\item
  Index: \texttt{\{\ video\_id,\ title,\ tags,\ transcript\_excerpt,\ view\_count,\ created\_at\ \}}.
\item
  Full-text + vector embedding for semantic search.
\end{itemize}

\textbf{Blob Storage --- S3-compatible object store (AWS S3 / GCS):}

\begin{itemize}
\tightlist
\item
  Raw uploads in one bucket; transcoded segments in another bucket per region.
\item
  Lifecycle policy: raw deleted after 30 days if status=ready.
\end{itemize}

\subsection{5. High-Level Architecture}\label{5-high-level-architecture}

\setcodelang{}

\begin{verbatim}
                       ┌──────────────┐
  Creator  ──upload──► │  Upload API   │
                       │  (chunked,    │
                       │  multipart)   │
                       └──────┬───────┘
                              │ raw video → S3 (raw bucket)
                              │
                       ┌──────▼────────────────────┐
                       │  Message Queue (Kafka)     │
                       │  topic: video.uploaded     │
                       └──────┬────────────────────┘
                              │
             ┌────────────────▼───────────────────────┐
             │        Transcoding Workers (fleet)      │
             │  FFmpeg containers (k8s jobs)           │
             │  output: HLS segments per resolution    │
             └────────────────┬───────────────────────┘
                              │ segments → S3 (cdn-origin bucket)
                              │ manifest → S3
                              │ update DB status = ready
                              ▼
                       ┌──────────────┐
                       │  Metadata DB  │      ◄── Video API (reads)
                       │  (Postgres/   │
                       │   Spanner)    │
                       └──────────────┘

  Viewer ──request──► │  Video API   │ ──► Metadata DB (manifest_url)
                       └──────┬───────┘
                              │  redirect to CDN
                              ▼
                       ┌──────────────────────────────────┐
                       │   CDN (Cloudfront / Fastly)      │
                       │   PoPs in 50+ regions            │
                       │   Cache: HLS manifests + segments│
                       └──────┬───────────────────────────┘
                              │ cache miss
                              ▼
                       ┌──────────────┐
                       │  S3 Origin   │
                       │  (segments)  │
                       └──────────────┘

  ┌──────────────────────────────────────────┐
  │  Recommendation Service                  │
  │  Kafka consumer → Feature store (Redis)  │
  │  ML model (Two-Tower / Collaborative     │
  │  Filtering) → pre-computed feed in Redis │
  └──────────────────────────────────────────┘
\end{verbatim}

\subsection{6. Deep Dives}\label{6-deep-dives}

\subsubsection{6.1 Upload Pipeline \& Transcoding}\label{61-upload-pipeline--transcoding}

\textbf{Chunked Upload:}

\begin{enumerate}
\def\labelenumi{\arabic{enumi}.}
\tightlist
\item
  Client calls \texttt{/initiate-upload} \(\rightarrow\) receives a pre-signed S3 multipart upload URL.
\item
  Client uploads 5--10 MB chunks in parallel (S3 multipart supports up to 10K parts).
\item
  Client calls \texttt{/complete} \(\rightarrow\) API assembles the multipart object in S3.
\item
  API publishes \texttt{\{\ video\_id,\ s3\_key\ \}} to Kafka topic \texttt{video.uploaded}.
\end{enumerate}

\textbf{Transcoding Workers:}

\begin{itemize}
\tightlist
\item
  Kubernetes Jobs pulled from Kafka; each worker pulls raw video, runs FFmpeg.
\item
  \textbf{Parallel transcoding:} one worker per resolution, or one worker per video with all resolutions in parallel child processes.
\item
  Output: HLS (\texttt{.m3u8} master manifest + \texttt{.ts} segment files per resolution, default 6 s segments).
\item
  Also output a DASH \texttt{MPD} manifest for Chrome/Android compatibility.
\item
  Store thumbnails (extracted at T=5s, 30s, 60s) in S3.
\item
  Update \texttt{video\_renditions} table and set \texttt{status\ =\ \textquotesingle{}ready\textquotesingle{}}.
\end{itemize}

\textbf{Resolution ladder (YouTube standard):}

\setcodelang{}

\begin{verbatim}
2160p (4K)   → bitrate ~20 Mbps
1080p        → ~8 Mbps
720p         → ~5 Mbps
480p         → ~2.5 Mbps
360p         → ~1 Mbps
240p         → ~0.5 Mbps
\end{verbatim}

\subsubsection{6.2 CDN Distribution}\label{62-cdn-distribution}

\begin{itemize}
\tightlist
\item
  S3 origin sits behind CDN (e.g., CloudFront with 400+ PoPs).
\item
  HLS segments are tiny (\textless{} 1 MB each, 6 s) --- extremely cache-friendly.
\item
  \textbf{Cache TTL:} segments are immutable \(\rightarrow\) TTL = \(\infty\) (or 30 days). Manifests refreshed every 60 s.
\item
  \textbf{Geo-routing:} DNS returns nearest PoP IP via Anycast/Route 53 Latency Routing.
\item
  \textbf{Cache hit ratio target:} \textgreater{} 95\% for popular videos. Long-tail videos miss to origin (cheap S3 GET).
\end{itemize}

\subsubsection{6.3 Adaptive Bitrate Streaming --- HLS vs DASH}\label{63-adaptive-bitrate-streaming--hls-vs-dash}

\begin{longtable}[]{@{}
  >{\raggedright\arraybackslash}p{(\columnwidth - 4\tabcolsep) * \real{0.1687}}
  >{\raggedright\arraybackslash}p{(\columnwidth - 4\tabcolsep) * \real{0.4578}}
  >{\raggedright\arraybackslash}p{(\columnwidth - 4\tabcolsep) * \real{0.3735}}@{}}
\toprule\noalign{}
\rowcolor{acc!16}
\begin{minipage}[b]{\linewidth}\raggedright
Feature
\end{minipage} & \begin{minipage}[b]{\linewidth}\raggedright
HLS (Apple)
\end{minipage} & \begin{minipage}[b]{\linewidth}\raggedright
DASH (MPEG)
\end{minipage} \\
\midrule\noalign{}
\endhead
\bottomrule\noalign{}
\endlastfoot
\textbf{Standard} & Apple proprietary \(\rightarrow\) RFC 8216 & MPEG open standard \\
\rowcolor{zebra}
\textbf{Container} & MPEG-TS or fMP4 & fMP4 (CMAF) \\
\textbf{iOS/Safari} & Native support & Requires MSE (not on older iOS) \\
\rowcolor{zebra}
\textbf{Android/Chrome} & Via MSE & Native support \\
\textbf{DRM} & FairPlay & Widevine / PlayReady \\
\rowcolor{zebra}
\textbf{Latency (live)} & 6--30 s (LLHLS: 1--3 s) & 2--8 s (CMAF chunked) \\
\textbf{Adoption} & YouTube, Netflix (backup) & Netflix (primary), Disney+ \\
\end{longtable}

\textbf{Key decision: Serve both HLS and DASH manifests from the same fMP4 segments (CMAF).} This halves storage vs maintaining two separate segment sets.

\subsubsection{6.4 Storage Tiers}\label{64-storage-tiers}

\setcodelang{}

\begin{verbatim}
Hot   (< 7 days old)   → S3 Standard          (fast GET, higher cost)
Warm  (7–90 days)      → S3 Infrequent Access (30% cheaper)
Cold  (> 90 days, < 1M views)  → S3 Glacier Instant Retrieval
Archive (> 1yr, < 10K views)   → S3 Glacier Deep Archive (pennies/GB)
\end{verbatim}

\subsubsection{6.5 Recommendations (brief)}\label{65-recommendations-brief}

\begin{itemize}
\tightlist
\item
  \textbf{Candidate generation:} Two-Tower model; user embedding \(\times\) video embedding \(\rightarrow\) top-K ANN (FAISS / ScaNN).
\item
  \textbf{Ranking:} LightGBM or DNN re-ranker on features: CTR, watch time, freshness, diversity.
\item
  \textbf{Serving:} Pre-compute homepage feed into Redis (TTL 30 min); real-time session signals update weights.
\item
  \textbf{Cold start:} New videos bootstrapped with category/tag-based rules.
\end{itemize}

\subsection{7. Bottlenecks \& Scaling}\label{7-bottlenecks--scaling}

\begin{longtable}[]{@{}
  >{\raggedright\arraybackslash}p{(\columnwidth - 2\tabcolsep) * \real{0.3866}}
  >{\raggedright\arraybackslash}p{(\columnwidth - 2\tabcolsep) * \real{0.6134}}@{}}
\toprule\noalign{}
\rowcolor{acc!16}
\begin{minipage}[b]{\linewidth}\raggedright
Bottleneck
\end{minipage} & \begin{minipage}[b]{\linewidth}\raggedright
Mitigation
\end{minipage} \\
\midrule\noalign{}
\endhead
\bottomrule\noalign{}
\endlastfoot
Transcoding backlog & Auto-scale k8s transcoding pods on Kafka consumer lag metric \\
\rowcolor{zebra}
Metadata DB hot rows (view counts) & Counter service with Redis; async batch-write to DB \\
CDN cache misses for new videos & Pre-warm CDN by pushing manifest URL to CDN after transcoding \\
\rowcolor{zebra}
S3 request throttling & Randomize S3 key prefixes (avoid date-prefix hotspot) \\
Recommendation latency & Pre-compute + cache feeds; async refresh \\
\rowcolor{zebra}
Single-region failure & Multi-region S3 replication; CDN absorbs reads; DB cross-region replica \\
\end{longtable}

\subsection{8. Trade-offs Summary}\label{8-trade-offs-summary}

\begin{longtable}[]{@{}
  >{\raggedright\arraybackslash}p{(\columnwidth - 6\tabcolsep) * \real{0.2895}}
  >{\raggedright\arraybackslash}p{(\columnwidth - 6\tabcolsep) * \real{0.2018}}
  >{\raggedright\arraybackslash}p{(\columnwidth - 6\tabcolsep) * \real{0.1930}}
  >{\raggedright\arraybackslash}p{(\columnwidth - 6\tabcolsep) * \real{0.3158}}@{}}
\toprule\noalign{}
\rowcolor{acc!16}
\begin{minipage}[b]{\linewidth}\raggedright
Decision
\end{minipage} & \begin{minipage}[b]{\linewidth}\raggedright
Option A
\end{minipage} & \begin{minipage}[b]{\linewidth}\raggedright
Option B
\end{minipage} & \begin{minipage}[b]{\linewidth}\raggedright
\textbf{Choice}
\end{minipage} \\
\midrule\noalign{}
\endhead
\bottomrule\noalign{}
\endlastfoot
Transcoding: per-video vs sharded & Single big VM per video & K8s pod per resolution & \textbf{K8s pods} (parallelism, cost) \\
\rowcolor{zebra}
Metadata DB & MySQL sharded & Spanner / CockroachDB & \textbf{Spanner} if global consistency needed \\
Streaming format & HLS-only & HLS + DASH (CMAF) & \textbf{Both} (broadest device coverage) \\
\rowcolor{zebra}
Recommendation serving & Real-time inference & Pre-computed feed & \textbf{Pre-computed} (latency wins) \\
Storage tiering & All S3 Standard & Tiered (hot/warm/cold) & \textbf{Tiered} (10\(\times\) cost reduction) \\
\end{longtable}

\textbf{Interview takeaways:}

\begin{itemize}
\tightlist
\item
  \textbf{The upload path and the playback path are entirely decoupled} --- async Kafka pipeline.
\item
  \textbf{CDN cache hit ratio is the single most important scaling lever} for a video platform.
\item
  \textbf{CMAF lets you serve HLS and DASH from the same segments} --- always mention this.
\item
  Transcoding is CPU-bound; scale it independently with Kafka consumer-lag autoscaling.
\end{itemize}

\section{2. Design a Ride-Hailing Service (Uber / Lyft)}\label{2-design-a-ride-hailing-service-uber--lyft}

\subsection{1. Requirements}\label{1-requirements-1}

\textbf{Clarifying questions:}

\begin{itemize}
\tightlist
\item
  Real-time ride-matching, or scheduled rides too?
\item
  Single city or global deployment?
\item
  Do we need the driver payout/financial system?
\item
  Surge pricing --- dynamic or simple multipliers?
\item
  Bike/scooter/autonomous vehicles, or just cars?
\end{itemize}

\textbf{Functional Requirements:}

\begin{itemize}
\tightlist
\item
  Riders request a ride from location A to B.
\item
  System matches rider to nearby available driver in real time.
\item
  Driver and rider track each other on a live map during the trip.
\item
  Compute fare estimate upfront; charge on trip completion.
\item
  Surge pricing during high-demand periods.
\item
  Driver/rider rating after trip.
\end{itemize}

\textbf{Non-Functional Requirements:}

\begin{itemize}
\tightlist
\item
  Matching latency \textless{} 3 s p99 from rider request to driver offer sent.
\item
  Location update ingestion: 1--5 s refresh from every active driver.
\item
  Availability: 99.99\% (5 min/month downtime is unacceptable during peak hours).
\item
  Consistency: eventual for location data; strong for trip state (no double-booking a driver).
\item
  Global scale: 5M+ concurrent drivers; 10M+ concurrent riders during peak.
\end{itemize}

\subsection{2. Capacity Estimation}\label{2-capacity-estimation-1}

\setcodelang{Python}

\begin{Shaded}
\begin{Highlighting}[]
\NormalTok{Active drivers (peak }\KeywordTok{global}\NormalTok{)   }\OperatorTok{=} \DecValTok{5}\NormalTok{ million}
\NormalTok{Location update interval       }\OperatorTok{=} \DecValTok{4}\NormalTok{ s}
\NormalTok{Location update QPS            }\OperatorTok{=} \DecValTok{5}\NormalTok{,}\DecValTok{000}\NormalTok{,}\DecValTok{000} \OperatorTok{/} \DecValTok{4} \OperatorTok{=} \DecValTok{1}\NormalTok{,}\DecValTok{250}\NormalTok{,}\DecValTok{000}\NormalTok{ writes}\OperatorTok{/}\NormalTok{sec}

\NormalTok{Active riders requesting rides }\OperatorTok{=} \DecValTok{500}\NormalTok{,}\DecValTok{000}\NormalTok{ concurrent}
\NormalTok{Ride requests QPS              }\OperatorTok{=} \DecValTok{500}\NormalTok{,}\DecValTok{000} \OperatorTok{/} \DecValTok{10}\ErrorTok{s}\NormalTok{ avg session }\OperatorTok{=} \DecValTok{50}\NormalTok{,}\DecValTok{000}\NormalTok{ QPS}

\NormalTok{Trips per day                  }\OperatorTok{=} \DecValTok{25}\NormalTok{ million (Uber }\OperatorTok{\textasciitilde{}}\DecValTok{20}\ErrorTok{M}\OperatorTok{/}\NormalTok{day)}
\NormalTok{Trip data per trip             }\OperatorTok{=} \OperatorTok{\textasciitilde{}}\DecValTok{2}\NormalTok{ KB metadata }\OperatorTok{+}\NormalTok{ GPS trace }\OperatorTok{\textasciitilde{}}\DecValTok{50}\NormalTok{ KB}
\NormalTok{Storage}\OperatorTok{/}\NormalTok{yr                     }\OperatorTok{=} \DecValTok{25}\ErrorTok{M}\NormalTok{ × }\DecValTok{52}\NormalTok{ KB × }\DecValTok{365} \OperatorTok{=} \OperatorTok{\textasciitilde{}}\DecValTok{475}\NormalTok{ TB}\OperatorTok{/}\NormalTok{yr}

\NormalTok{Location store size (}\KeywordTok{in}\OperatorTok{{-}}\NormalTok{memory):}
  \DecValTok{5}\ErrorTok{M}\NormalTok{ drivers × (lat, lon, heading, timestamp, driver\_id) }\OperatorTok{=} \OperatorTok{\textasciitilde{}}\DecValTok{200} \BuiltInTok{bytes}\NormalTok{ each}
  \OperatorTok{=} \DecValTok{5}\ErrorTok{M}\NormalTok{ × }\DecValTok{200}\NormalTok{ B }\OperatorTok{=} \DecValTok{1}\NormalTok{ GB — fits }\KeywordTok{in}\NormalTok{ a large Redis cluster}
\end{Highlighting}
\end{Shaded}

\subsection{3. API Design}\label{3-api-design-1}

\setcodelang{Python}

\begin{Shaded}
\begin{Highlighting}[]
\CommentTok{\# Rider}
\NormalTok{POST   }\OperatorTok{/}\NormalTok{v1}\OperatorTok{/}\NormalTok{rides}\OperatorTok{/}\NormalTok{estimate}
\NormalTok{       Body: \{ pickup\_lat, pickup\_lon, dest\_lat, dest\_lon, ride\_type \}}
\NormalTok{       Returns: \{ eta\_s, fare\_estimate, surge\_multiplier \}}

\NormalTok{POST   }\OperatorTok{/}\NormalTok{v1}\OperatorTok{/}\NormalTok{rides}\OperatorTok{/}\NormalTok{request}
\NormalTok{       Body: \{ pickup\_lat, pickup\_lon, dest\_lat, dest\_lon, ride\_type \}}
\NormalTok{       Returns: \{ ride\_id, status: }\StringTok{"searching"}\NormalTok{ \}}

\NormalTok{GET    }\OperatorTok{/}\NormalTok{v1}\OperatorTok{/}\NormalTok{rides}\OperatorTok{/}\NormalTok{\{ride\_id\}}
\NormalTok{       Returns: \{ status, driver\_id, driver\_eta\_s, driver\_location \}}

\NormalTok{DELETE }\OperatorTok{/}\NormalTok{v1}\OperatorTok{/}\NormalTok{rides}\OperatorTok{/}\NormalTok{\{ride\_id\}      (cancel)}

\CommentTok{\# Driver}
\NormalTok{PUT    }\OperatorTok{/}\NormalTok{v1}\OperatorTok{/}\NormalTok{drivers}\OperatorTok{/}\NormalTok{\{driver\_id\}}\OperatorTok{/}\NormalTok{location}
\NormalTok{       Body: \{ lat, lon, heading, speed, timestamp \}}

\NormalTok{PUT    }\OperatorTok{/}\NormalTok{v1}\OperatorTok{/}\NormalTok{drivers}\OperatorTok{/}\NormalTok{\{driver\_id\}}\OperatorTok{/}\NormalTok{status}
\NormalTok{       Body: \{ status: }\StringTok{"available"} \OperatorTok{|} \StringTok{"busy"} \OperatorTok{|} \StringTok{"offline"}\NormalTok{ \}}

\NormalTok{GET    }\OperatorTok{/}\NormalTok{v1}\OperatorTok{/}\NormalTok{drivers}\OperatorTok{/}\NormalTok{\{driver\_id\}}\OperatorTok{/}\NormalTok{ride}\OperatorTok{{-}}\NormalTok{offers   (}\BuiltInTok{long}\OperatorTok{{-}}\NormalTok{poll }\KeywordTok{or}\NormalTok{ WebSocket push)}
\NormalTok{POST   }\OperatorTok{/}\NormalTok{v1}\OperatorTok{/}\NormalTok{ride}\OperatorTok{{-}}\NormalTok{offers}\OperatorTok{/}\NormalTok{\{offer\_id\}}\OperatorTok{/}\NormalTok{accept}
\NormalTok{POST   }\OperatorTok{/}\NormalTok{v1}\OperatorTok{/}\NormalTok{ride}\OperatorTok{{-}}\NormalTok{offers}\OperatorTok{/}\NormalTok{\{offer\_id\}}\OperatorTok{/}\NormalTok{reject}

\CommentTok{\# Shared real{-}time (WebSocket)}
\NormalTok{WS }\OperatorTok{/}\NormalTok{v1}\OperatorTok{/}\NormalTok{ws}\OperatorTok{/}\NormalTok{ride}\OperatorTok{/}\NormalTok{\{ride\_id\}   (bidirectional: location pings, status events)}
\end{Highlighting}
\end{Shaded}

\subsection{4. Data Model}\label{4-data-model-1}

\textbf{Drivers (PostgreSQL):}

\setcodelang{SQL}

\begin{Shaded}
\begin{Highlighting}[]
\NormalTok{drivers (}
\NormalTok{  driver\_id       BIGINT }\KeywordTok{PRIMARY} \KeywordTok{KEY}\NormalTok{,}
\NormalTok{  name            TEXT,}
\NormalTok{  vehicle\_id      BIGINT,}
\NormalTok{  rating          }\DataTypeTok{NUMERIC}\NormalTok{(}\DecValTok{3}\NormalTok{,}\DecValTok{2}\NormalTok{),}
\NormalTok{  current\_status  ENUM(}\StringTok{\textquotesingle{}available\textquotesingle{}}\NormalTok{,}\StringTok{\textquotesingle{}busy\textquotesingle{}}\NormalTok{,}\StringTok{\textquotesingle{}offline\textquotesingle{}}\NormalTok{),}
\NormalTok{  city\_id         }\DataTypeTok{INT}
\NormalTok{)}
\end{Highlighting}
\end{Shaded}

\textbf{Rides (PostgreSQL sharded by ride\_id):}

\setcodelang{SQL}

\begin{Shaded}
\begin{Highlighting}[]
\NormalTok{rides (}
\NormalTok{  ride\_id         UUID }\KeywordTok{PRIMARY} \KeywordTok{KEY}\NormalTok{,}
\NormalTok{  rider\_id        BIGINT,}
\NormalTok{  driver\_id       BIGINT,}
\NormalTok{  status          ENUM(}\StringTok{\textquotesingle{}requested\textquotesingle{}}\NormalTok{,}\StringTok{\textquotesingle{}searching\textquotesingle{}}\NormalTok{,}\StringTok{\textquotesingle{}driver\_assigned\textquotesingle{}}\NormalTok{,}\StringTok{\textquotesingle{}in\_progress\textquotesingle{}}\NormalTok{,}\StringTok{\textquotesingle{}completed\textquotesingle{}}\NormalTok{,}\StringTok{\textquotesingle{}cancelled\textquotesingle{}}\NormalTok{),}
\NormalTok{  pickup\_lat      }\DataTypeTok{DOUBLE}\NormalTok{,}
\NormalTok{  pickup\_lon      }\DataTypeTok{DOUBLE}\NormalTok{,}
\NormalTok{  dest\_lat        }\DataTypeTok{DOUBLE}\NormalTok{,}
\NormalTok{  dest\_lon        }\DataTypeTok{DOUBLE}\NormalTok{,}
\NormalTok{  requested\_at    TIMESTAMPTZ,}
\NormalTok{  matched\_at      TIMESTAMPTZ,}
\NormalTok{  completed\_at    TIMESTAMPTZ,}
\NormalTok{  fare\_estimate   }\DataTypeTok{NUMERIC}\NormalTok{(}\DecValTok{10}\NormalTok{,}\DecValTok{2}\NormalTok{),}
\NormalTok{  final\_fare      }\DataTypeTok{NUMERIC}\NormalTok{(}\DecValTok{10}\NormalTok{,}\DecValTok{2}\NormalTok{),}
\NormalTok{  surge\_mult      }\DataTypeTok{NUMERIC}\NormalTok{(}\DecValTok{4}\NormalTok{,}\DecValTok{2}\NormalTok{)}
\NormalTok{)}
\end{Highlighting}
\end{Shaded}

\textbf{Driver Locations (Redis / in-memory geo store):}

\setcodelang{Python}

\begin{Shaded}
\begin{Highlighting}[]
\NormalTok{Key:   DRIVER\_LOC:\{driver\_id\}}
\NormalTok{Value: \{ lat, lon, heading, updated\_at \}}
\NormalTok{TTL:   }\DecValTok{30}\NormalTok{ s (stale }\ControlFlowTok{if}\NormalTok{ driver stops sending heartbeats)}

\NormalTok{Geo index: Redis GEOADD city:\{city\_id\}:drivers \{lon\} \{lat\} \{driver\_id\}}
\NormalTok{           Redis GEORADIUS }\ControlFlowTok{for}\NormalTok{ nearest}\OperatorTok{{-}}\NormalTok{N lookup}
\end{Highlighting}
\end{Shaded}

\textbf{Trip GPS traces (Cassandra or TimescaleDB):}

\setcodelang{}

\begin{verbatim}
trip_locations (
  ride_id     UUID,
  ts          TIMESTAMPTZ,
  lat         DOUBLE,
  lon         DOUBLE,
  PRIMARY KEY (ride_id, ts)
) — append-only, time-series access pattern
\end{verbatim}

\subsection{5. High-Level Architecture}\label{5-high-level-architecture-1}

\setcodelang{}

\begin{verbatim}
  Rider App                          Driver App
      │                                  │
      │ HTTPS                            │ HTTPS
      ▼                                  ▼
  ┌─────────────────────────────────────────────┐
  │          API Gateway / Load Balancer         │
  └──────┬──────────────────────────┬───────────┘
         │                          │
  ┌──────▼──────┐            ┌──────▼──────────┐
  │  Ride        │            │  Location        │
  │  Service     │            │  Service         │
  │  (stateless) │            │  (stateless)     │
  └──────┬───────┘            └──────┬───────────┘
         │                           │
         │                    ┌──────▼────────────────────┐
         │                    │  Location Store            │
         │                    │  Redis Cluster (Geo)       │
         │                    │  5M driver locations       │
         │                    └───────────────────────────┘
         │
  ┌──────▼────────────────────────────────────┐
  │         Matching Service                   │
  │  - Query Location Store for nearby drivers │
  │  - Score + rank candidates                 │
  │  - Send offer to best driver (WebSocket)   │
  │  - Timeout → next best driver              │
  └──────┬────────────────────────────────────┘
         │
  ┌──────▼──────────────┐  ┌──────────────────────┐
  │   Rides DB           │  │  Surge Pricing        │
  │   (Postgres sharded) │  │  Service              │
  │                      │  │  (demand/supply ratio │
  │                      │  │   per geohash cell)   │
  └──────────────────────┘  └──────────────────────┘
         │
  ┌──────▼──────────────┐
  │  Notification        │
  │  Service             │
  │  (FCM / APNs push)   │
  └──────────────────────┘

  WebSocket Gateway (separate fleet, stateful):
  ┌──────────────────────────────────────────┐
  │  Rider WS connection ◄──► Ride Service   │
  │  Driver WS connection ◄──► Ride Service  │
  │  (publish/subscribe via Redis Pub/Sub)   │
  └──────────────────────────────────────────┘
\end{verbatim}

\subsection{6. Deep Dives}\label{6-deep-dives-1}

\subsubsection{6.1 Geospatial Indexing --- Geohash vs Quadtree vs S2}\label{61-geospatial-indexing--geohash-vs-quadtree-vs-s2}

\begin{longtable}[]{@{}
  >{\raggedright\arraybackslash}p{(\columnwidth - 6\tabcolsep) * \real{0.1574}}
  >{\raggedright\arraybackslash}p{(\columnwidth - 6\tabcolsep) * \real{0.2963}}
  >{\raggedright\arraybackslash}p{(\columnwidth - 6\tabcolsep) * \real{0.2315}}
  >{\raggedright\arraybackslash}p{(\columnwidth - 6\tabcolsep) * \real{0.3148}}@{}}
\toprule\noalign{}
\rowcolor{acc!16}
\begin{minipage}[b]{\linewidth}\raggedright
Feature
\end{minipage} & \begin{minipage}[b]{\linewidth}\raggedright
Geohash
\end{minipage} & \begin{minipage}[b]{\linewidth}\raggedright
Quadtree
\end{minipage} & \begin{minipage}[b]{\linewidth}\raggedright
Google S2
\end{minipage} \\
\midrule\noalign{}
\endhead
\bottomrule\noalign{}
\endlastfoot
\textbf{Structure} & Base-32 string, hierarchical & Tree of rectangular cells & Hilbert-curve cells on sphere \\
\rowcolor{zebra}
\textbf{Precision control} & String length (1--12 chars) & Tree depth & Level (0--30) \\
\textbf{Query type} & String prefix match & Tree traversal & Cell union + covering \\
\rowcolor{zebra}
\textbf{Edge distortion} & Yes (poles, antimeridian) & Yes (aspect ratio varies) & Minimal (spherical geometry) \\
\textbf{Neighbor lookup} & 8 neighbors (manual computation) & Tree sibling traversal & S2::GetEdgeNeighbors \\
\rowcolor{zebra}
\textbf{Used by} & MongoDB, Elasticsearch, Redis & Older game engines & Uber (H3), Google Maps, Foursquare \\
\textbf{Complexity} & Simple & Moderate & Higher (library dependency) \\
\end{longtable}

\textbf{Key decision: Use Geohash (precision level 6 \(\approx\) 1.2 km cells) for driver lookup.}

\begin{itemize}
\tightlist
\item
  Each driver is stored in Redis with \texttt{GEOADD} (which internally uses Geohash).
\item
  On ride request: \texttt{GEORADIUS\ pickup\_lat\ pickup\_lon\ 3km} \(\rightarrow\) list of driver IDs \(\rightarrow\) filter by \texttt{status=available}.
\item
  \textbf{Why not quadtree?} Redis GEODATA gives us geohash for free; quadtree requires custom in-memory structure.
\item
  \textbf{Why not S2?} Overkill for matching; S2 shines in complex polygon coverage (e.g., Uber H3 for analytics).
\end{itemize}

\subsubsection{6.2 Driver-Rider Matching}\label{62-driver-rider-matching}

\textbf{Algorithm:}

\begin{enumerate}
\def\labelenumi{\arabic{enumi}.}
\tightlist
\item
  \texttt{GEORADIUS} call returns N available drivers within 3 km (expand to 5 km if \textless{} 3 results).
\item
  Score each driver:

  \setcodelang{Python}

\begin{Shaded}
\begin{Highlighting}[]
\NormalTok{score }\OperatorTok{=}\NormalTok{ w1 × (}\DecValTok{1} \OperatorTok{/}\NormalTok{ ETA\_to\_pickup) }\OperatorTok{+}\NormalTok{ w2 × driver\_rating }\OperatorTok{+}\NormalTok{ w3 × acceptance\_rate}
\end{Highlighting}
\end{Shaded}
\item
  Send offer to top-scored driver. Wait 10 s.
\item
  If rejected/timeout \(\rightarrow\) send to next driver. Max 5 tries before returning "no drivers available."
\item
  On acceptance \(\rightarrow\) atomically set \texttt{rides.driver\_id}, \texttt{rides.status\ =\ driver\_assigned}, \texttt{drivers.current\_status\ =\ busy} (Postgres transaction).
\end{enumerate}

\textbf{Preventing double-assignment:} Use Postgres \texttt{SELECT\ FOR\ UPDATE} or Redis distributed lock (\texttt{SET\ driver:\{id\}:lock\ NX\ EX\ 15}) before assigning.

\subsubsection{6.3 Real-Time Location Updates}\label{63-real-time-location-updates}

\begin{itemize}
\tightlist
\item
  Drivers send GPS ping every 4 s via HTTP/2 (or WebSocket for persistent connection).
\item
  \textbf{Location Service} writes to Redis \texttt{GEOADD} (O(log N)) --- 1.25M writes/sec requires Redis Cluster with 10--20 shards.
\item
  \textbf{Sharding:} partition by \texttt{city\_id} or geohash prefix (level 2 \(\approx\) 630 km cells \(\rightarrow\) \textasciitilde100 world cells \(\rightarrow\) manageable shards).
\item
  Driver location also published to \textbf{Kafka} topic \texttt{driver.location} for:

  \begin{itemize}
  \tightlist
  \item
    Live map updates pushed to rider via WebSocket.
  \item
    Surge pricing calculation.
  \item
    ETA recalculation.
  \item
    Trip trace storage (Cassandra consumer).
  \end{itemize}
\end{itemize}

\subsubsection{6.4 Surge Pricing}\label{64-surge-pricing}

\setcodelang{Python}

\begin{Shaded}
\begin{Highlighting}[]
\NormalTok{demand\_cells[geohash6] }\OperatorTok{=}\NormalTok{ count of ride requests }\KeywordTok{in}\NormalTok{ last }\DecValTok{5} \BuiltInTok{min}
\NormalTok{supply\_cells[geohash6] }\OperatorTok{=}\NormalTok{ count of available drivers }\KeywordTok{in}\NormalTok{ cell}

\NormalTok{surge\_multiplier }\OperatorTok{=} \BuiltInTok{max}\NormalTok{(}\FloatTok{1.0}\NormalTok{, demand }\OperatorTok{/}\NormalTok{ (supply × target\_utilization\_factor))}
\end{Highlighting}
\end{Shaded}

\begin{itemize}
\tightlist
\item
  Computed every 30 s by Surge Service consuming Kafka streams.
\item
  Stored in Redis: \texttt{SURGE:\{geohash6\}} \(\rightarrow\) multiplier (float).
\item
  Ride estimate reads from Redis before quoting price.
\item
  \textbf{UI UX:} show surge visually (heatmap) and require explicit rider confirmation if multiplier \textgreater{} 1.5\(\times\).
\end{itemize}

\subsubsection{6.5 Trip Lifecycle --- State Machine}\label{65-trip-lifecycle--state-machine}

\setcodelang{}

\begin{verbatim}
REQUESTED ──► SEARCHING ──► DRIVER_ASSIGNED ──► DRIVER_EN_ROUTE
                                                      │
                                                      ▼
                                               ARRIVED_AT_PICKUP
                                                      │
                                                      ▼
                                                IN_PROGRESS
                                                      │
                                                      ▼
                                               COMPLETED / CANCELLED
\end{verbatim}

State transitions are persisted in Postgres with timestamps. Events published to Kafka on each transition. Kafka consumers handle: notifications, billing trigger (on COMPLETED), analytics.

\subsubsection{6.6 ETA Calculation}\label{66-eta-calculation}

\begin{itemize}
\tightlist
\item
  \textbf{Routing engine:} OSRM or Google Maps Platform API call on driver assignment.
\item
  \textbf{Real-time recalculation:} every 30 s while driver is en route, re-query routing engine with current driver location.
\item
  \textbf{Pre-computed ETA surfaces:} city road graph stored in-memory (contraction hierarchies) for sub-100ms routing. Uber uses their own in-house routing (H3-based graph).
\end{itemize}

\subsection{7. Bottlenecks \& Scaling}\label{7-bottlenecks--scaling-1}

\begin{longtable}[]{@{}
  >{\raggedright\arraybackslash}p{(\columnwidth - 2\tabcolsep) * \real{0.3652}}
  >{\raggedright\arraybackslash}p{(\columnwidth - 2\tabcolsep) * \real{0.6348}}@{}}
\toprule\noalign{}
\rowcolor{acc!16}
\begin{minipage}[b]{\linewidth}\raggedright
Bottleneck
\end{minipage} & \begin{minipage}[b]{\linewidth}\raggedright
Mitigation
\end{minipage} \\
\midrule\noalign{}
\endhead
\bottomrule\noalign{}
\endlastfoot
Location write storm & Redis Cluster, shard by city/geohash; accept at-most-once semantics \\
\rowcolor{zebra}
Matching hot cities (NYC peak) & City-level matching service instances; limit to local Redis shard \\
Rides DB write contention & Shard by \texttt{city\_id}; use UUIDs to avoid hot auto-increment keys \\
\rowcolor{zebra}
WebSocket connection scaling & Stateful WS servers behind L4 LB; use Redis Pub/Sub as broker \\
Driver double-booking & Redis NX lock on \texttt{driver\_id} during match; release on timeout \\
\rowcolor{zebra}
Surge pricing staleness & Short TTL (30 s) + frequent recalculation; eventual consistency OK \\
\end{longtable}

\subsection{8. Trade-offs Summary}\label{8-trade-offs-summary-1}

\begin{longtable}[]{@{}
  >{\raggedright\arraybackslash}p{(\columnwidth - 6\tabcolsep) * \real{0.2617}}
  >{\raggedright\arraybackslash}p{(\columnwidth - 6\tabcolsep) * \real{0.2150}}
  >{\raggedright\arraybackslash}p{(\columnwidth - 6\tabcolsep) * \real{0.2056}}
  >{\raggedright\arraybackslash}p{(\columnwidth - 6\tabcolsep) * \real{0.3178}}@{}}
\toprule\noalign{}
\rowcolor{acc!16}
\begin{minipage}[b]{\linewidth}\raggedright
Decision
\end{minipage} & \begin{minipage}[b]{\linewidth}\raggedright
Option A
\end{minipage} & \begin{minipage}[b]{\linewidth}\raggedright
Option B
\end{minipage} & \begin{minipage}[b]{\linewidth}\raggedright
\textbf{Choice}
\end{minipage} \\
\midrule\noalign{}
\endhead
\bottomrule\noalign{}
\endlastfoot
Geospatial index & Geohash (Redis GEODATA) & Custom Quadtree & \textbf{Geohash (Redis)} \\
\rowcolor{zebra}
Driver location transport & HTTP polling (4 s) & Persistent WebSocket & \textbf{HTTP/2} (simpler, scales well) \\
Matching concurrency control & Postgres FOR UPDATE & Redis distributed lock & \textbf{Redis lock} (lower latency) \\
\rowcolor{zebra}
Location store & Redis in-memory & PostGIS & \textbf{Redis} (1M+ writes/sec) \\
Trip state persistence & Postgres & Cassandra & \textbf{Postgres} (ACID for state machines) \\
\rowcolor{zebra}
Routing engine & Google Maps API & Self-hosted OSRM & \textbf{OSRM} at scale (cost) \\
\end{longtable}

\textbf{Interview takeaways:}

\begin{itemize}
\tightlist
\item
  \textbf{The geospatial indexing approach is the core of this design} --- be ready to explain Geohash precision levels.
\item
  \textbf{Driver location is write-heavy, trip state is consistency-critical} --- use different stores for each.
\item
  \textbf{The matching loop is an iterative fan-out} --- must handle timeouts and fallbacks gracefully.
\item
  Surge pricing uses supply/demand ratio per geohash cell --- a simple but powerful formula.
\end{itemize}

\section{3. Design a Distributed Cache (like Redis/Memcached at Scale)}\label{3-design-a-distributed-cache-like-redismemcached-at-scale}

\subsection{1. Requirements}\label{1-requirements-2}

\textbf{Clarifying questions:}

\begin{itemize}
\tightlist
\item
  Key-value only, or also sorted sets, pub/sub, streams?
\item
  What are the SLAs on read/write latency?
\item
  Strong consistency needed, or eventual?
\item
  Target use case: session cache, DB cache, rate limiter, leaderboard?
\item
  What\textquotesingle s the expected working set size (total hot data)?
\end{itemize}

\textbf{Functional Requirements:}

\begin{itemize}
\tightlist
\item
  \texttt{GET\ key}, \texttt{SET\ key\ value\ {[}TTL{]}}, \texttt{DEL\ key}, \texttt{MGET}, \texttt{MSET}.
\item
  TTL-based expiration.
\item
  LRU eviction when memory is full.
\item
  Replication for high availability.
\item
  Horizontal scaling via consistent hashing.
\item
  Pub/Sub (optional advanced feature).
\end{itemize}

\textbf{Non-Functional Requirements:}

\begin{itemize}
\tightlist
\item
  Read latency \textless{} 1 ms p99.
\item
  Write latency \textless{} 2 ms p99.
\item
  Availability: 99.99\% (cache miss falls through to DB, but should be rare).
\item
  No data loss for critical caches (optional persistence via WAL/RDB).
\item
  Scale to terabytes of cached data across a cluster.
\end{itemize}

\subsection{2. Capacity Estimation}\label{2-capacity-estimation-2}

\setcodelang{}

\begin{verbatim}
Cached data size (working set)   = 10 TB
Node RAM (cache node)            = 256 GB
Nodes needed                     = 10 TB / 256 GB ≈ 40 cache nodes

Read QPS                         = 1,000,000 reads/sec (1M QPS)
Write QPS                        = 100,000 writes/sec (10% of reads)

Per-node QPS (40 nodes)          = 25,000 reads + 2,500 writes per node
Single Redis node can handle     ≈ 100K–200K ops/sec → 40 nodes has 4–8× headroom

Network bandwidth per node:
  25K reads × avg 1 KB value     = 25 MB/s per node  (well within 10 GbE NIC)

Memory per key overhead (Redis): ~70 bytes metadata + value size
  1B keys × 70 B overhead        = 70 GB overhead in cluster (significant at scale)
\end{verbatim}

\subsection{3. API Design}\label{3-api-design-2}

\textbf{Client-facing (wire protocol --- Redis-compatible RESP):}

\setcodelang{}

\begin{verbatim}
GET  <key>                   → value | (nil)
SET  <key> <value> [EX secs] → OK
DEL  <key> [key ...]         → (integer) count deleted
MGET <key> [key ...]         → array of values
MSET <key> <val> [key val]   → OK
TTL  <key>                   → seconds remaining
INCR / DECR <key>            → new integer value
EXPIRE <key> <secs>          → 1 (success) | 0 (no such key)
\end{verbatim}

\textbf{Admin / Internal API:}

\setcodelang{}

\begin{verbatim}
CLUSTER ADDNODE <host:port>  → join node to cluster
CLUSTER FAILOVER             → trigger manual leader promotion
INFO                         → stats (memory, QPS, hit_rate, evictions)
BGSAVE                       → async RDB snapshot
\end{verbatim}

\subsection{4. Data Model}\label{4-data-model-2}

\textbf{In-memory data structures:}

\setcodelang{}

\begin{verbatim}
Hash table: key → { value, type, TTL_expiry_ts, LRU_clock }

Value types:
  String → raw bytes (up to 512 MB)
  List   → doubly-linked list (small) or ziplist (< 128 elements)
  Hash   → hash table or ziplist
  Set    → hash table or intset
  ZSet   → skip list + hash table
\end{verbatim}

\textbf{Persistence (optional):}

\setcodelang{}

\begin{verbatim}
RDB snapshot: point-in-time binary dump, every N seconds or M writes
AOF (Append-Only File): write-ahead log, fsync policy:
  - always  → 1 write per command, safest, slowest
  - everysec → fsync every 1s, at most 1s of data loss
  - never   → OS decides, fastest, most data loss risk
\end{verbatim}

\textbf{Cluster topology metadata:}

\setcodelang{}

\begin{verbatim}
Cluster config (stored in every node):
  slot_assignments[0..16383] → node_id   (Redis uses 16384 hash slots)
  node_registry: { node_id, host, port, role: master|replica, slots }
\end{verbatim}

\subsection{5. High-Level Architecture}\label{5-high-level-architecture-2}

\setcodelang{}

\begin{verbatim}
  Client Apps (thousands of instances)
         │
  ┌──────▼──────────────────────────────────────────┐
  │         Client Library (smart client)            │
  │  - Consistent hash → route to correct node       │
  │  - Connection pool per node                      │
  │  - Retry on MOVED/ASK redirects                  │
  └──────┬──────────────────────────────────────────┘
         │  direct TCP (RESP protocol)
         ▼
  ┌──────────────────────────────────────────────────┐
  │            Cache Cluster (40 nodes)              │
  │                                                  │
  │  Shard 0       Shard 1       ...  Shard 19       │
  │  [Master M0]   [Master M1]        [Master M19]   │
  │  [Replica R0a] [Replica R1a]      [Replica R19a] │
  │  [Replica R0b] [Replica R1b]      [Replica R19b] │
  │                                                  │
  │  Each master owns 16384/20 ≈ 820 hash slots      │
  └──────────────────────────────────────────────────┘
         │
         │ cache miss
         ▼
  ┌──────────────────┐
  │  Backend DB       │
  │  (Postgres, etc.) │
  └──────────────────┘

  Cluster Coordination:
  ┌──────────────────────────────────────────────────┐
  │  Gossip Protocol (every node ↔ every other node) │
  │  - Heartbeat PING/PONG every 1 s                 │
  │  - Failure detection: missed 3× → PFAIL → FAIL   │
  │  - Automatic failover: replica elected master     │
  └──────────────────────────────────────────────────┘
\end{verbatim}

\subsection{6. Deep Dives}\label{6-deep-dives-2}

\subsubsection{6.1 Consistent Hashing}\label{61-consistent-hashing}

\textbf{Problem:} Naive \texttt{hash(key)\ \%\ N} rehashes ALL keys when N changes (add/remove node).

\textbf{Solution --- Consistent Hash Ring:}

\setcodelang{}

\begin{verbatim}
Virtual ring: 0 to 2^32
Each node placed at multiple positions (virtual nodes / vnodes)
  - 150 vnodes per physical node → even distribution
  - Rebalancing a new node: only transfer 1/N of keys (not all)

key → hash(key) → position on ring → clockwise next node = owner

                  0
                 /|\
           M3 ←  |  → M0
          /       |      \
        M2        |        M1
          \                /
            ─────────────
                2^32
\end{verbatim}

\textbf{Key decision:} 150--200 virtual nodes per physical node balances load even with heterogeneous node sizes (weight vnodes by RAM). With 40 nodes \(\times\) 150 vnodes = 6,000 points on ring.

\subsubsection{6.2 Replication}\label{62-replication}

\begin{itemize}
\tightlist
\item
  Each shard: 1 master + 2 replicas (RF=3).
\item
  \textbf{Async replication} (default Redis): master returns OK to client, replicates to replicas asynchronously. Risk: up to N ms of data loss on crash.
\item
  \textbf{Semi-sync (WAIT command):} \texttt{WAIT\ 1\ 100} --- wait for 1 replica to ack within 100ms. Better durability at slight latency cost.
\item
  Replicas serve \textbf{reads} to scale read-heavy workloads (with possible stale reads).
\item
  On master failure: replica with lowest replication lag is elected master (Raft-adjacent gossip vote).
\end{itemize}

\subsubsection{6.3 Eviction Policies (LRU and Beyond)}\label{63-eviction-policies-lru-and-beyond}

\begin{longtable}[]{@{}
  >{\raggedright\arraybackslash}p{(\columnwidth - 4\tabcolsep) * \real{0.1795}}
  >{\raggedright\arraybackslash}p{(\columnwidth - 4\tabcolsep) * \real{0.5128}}
  >{\raggedright\arraybackslash}p{(\columnwidth - 4\tabcolsep) * \real{0.3077}}@{}}
\toprule\noalign{}
\rowcolor{acc!16}
\begin{minipage}[b]{\linewidth}\raggedright
Policy
\end{minipage} & \begin{minipage}[b]{\linewidth}\raggedright
Description
\end{minipage} & \begin{minipage}[b]{\linewidth}\raggedright
Use case
\end{minipage} \\
\midrule\noalign{}
\endhead
\bottomrule\noalign{}
\endlastfoot
\textbf{noeviction} & Return error when memory full & Primary DB (never evict) \\
\rowcolor{zebra}
\textbf{allkeys-lru} & Evict globally least-recently-used key & General cache \\
\textbf{volatile-lru} & Evict LRU only among keys with TTL set & Mixed TTL/no-TTL data \\
\rowcolor{zebra}
\textbf{allkeys-lfu} & Evict least-frequently-used (Redis 4.0+) & Skewed access patterns \\
\textbf{volatile-ttl} & Evict key with shortest remaining TTL & Time-sensitive caching \\
\rowcolor{zebra}
\textbf{allkeys-random} & Evict random key & Uniform access, simplest \\
\end{longtable}

\textbf{Approximate LRU (Redis implementation):}

\begin{itemize}
\tightlist
\item
  Full LRU requires O(N) scan --- too slow.
\item
  Redis samples 5 (configurable) random keys, evicts the least-recently-used among them.
\item
  At sample=10, approximates true LRU with \textless{} 1\% error in most workloads.
\end{itemize}

\textbf{Key decision: Use \texttt{allkeys-lfu} for a general-purpose cache} --- LFU outperforms LRU when there are periodic scans that pollute the LRU queue with cold data.

\subsubsection{6.4 Cache Invalidation Strategies}\label{64-cache-invalidation-strategies}

\begin{longtable}[]{@{}
  >{\raggedright\arraybackslash}p{(\columnwidth - 6\tabcolsep) * \real{0.1990}}
  >{\raggedright\arraybackslash}p{(\columnwidth - 6\tabcolsep) * \real{0.4188}}
  >{\raggedright\arraybackslash}p{(\columnwidth - 6\tabcolsep) * \real{0.1672}}
  >{\raggedright\arraybackslash}p{(\columnwidth - 6\tabcolsep) * \real{0.2150}}@{}}
\toprule\noalign{}
\rowcolor{acc!16}
\begin{minipage}[b]{\linewidth}\raggedright
Strategy
\end{minipage} & \begin{minipage}[b]{\linewidth}\raggedright
Description
\end{minipage} & \begin{minipage}[b]{\linewidth}\raggedright
Pros
\end{minipage} & \begin{minipage}[b]{\linewidth}\raggedright
Cons
\end{minipage} \\
\midrule\noalign{}
\endhead
\bottomrule\noalign{}
\endlastfoot
\textbf{TTL expiry} & Set TTL on write; auto-expire & Simple & Stale window = TTL duration \\
\rowcolor{zebra}
\textbf{Write-through} & Write to cache + DB simultaneously & Always consistent & Write latency increased \\
\textbf{Write-behind (async)} & Write to cache, async flush to DB & Low write latency & Data loss risk \\
\rowcolor{zebra}
\textbf{Cache-aside} & App reads cache; miss \(\rightarrow\) read DB \(\rightarrow\) populate cache & DB is source of truth & Cache miss latency spike \\
\textbf{Read-through} & Cache library fetches from DB on miss automatically & Transparent to app & Library complexity \\
\rowcolor{zebra}
\textbf{Event-driven invalidation} & DB change \(\rightarrow\) CDC event \(\rightarrow\) DEL from cache & Real-time consistency & Infrastructure complexity \\
\end{longtable}

\textbf{Key decision: TTL + event-driven invalidation (CDC via Debezium/Kafka).} TTL covers eventual consistency; CDC invalidates on DB write for critical data.

\subsubsection{6.5 Hot Key Problem}\label{65-hot-key-problem}

\textbf{Problem:} A single key (e.g., a celebrity\textquotesingle s profile \texttt{user:12345}) gets 100K QPS --- overwhelming the one shard that owns it.

\textbf{Solutions:}

\begin{enumerate}
\def\labelenumi{\arabic{enumi}.}
\tightlist
\item
  \textbf{Key replication / local cache:} Client-side in-process cache (LRU, 100 MB per app server) for top-N hot keys. App caches key for 1 s locally.
\item
  \textbf{Scatter with suffix:} Store key as \texttt{user:12345:0} through \texttt{user:12345:9} (10 shards); reads pick random suffix, writes update all. Works for read-heavy hot keys.
\item
  \textbf{Read replicas:} Route reads to any of the replicas for that shard.
\item
  \textbf{Detect \& alert:} Monitor per-key QPS in the client library; auto-promote to local cache.
\end{enumerate}

\subsubsection{6.6 Cache Stampede (Thundering Herd)}\label{66-cache-stampede-thundering-herd}

\textbf{Problem:} Popular key expires \(\rightarrow\) 10K concurrent requests all miss \(\rightarrow\) all hit DB simultaneously \(\rightarrow\) DB overloaded.

\textbf{Solutions:}

\begin{enumerate}
\def\labelenumi{\arabic{enumi}.}
\tightlist
\item
  \textbf{Probabilistic early expiration (PER):} Recompute cache key slightly before TTL expires (during window \texttt{TTL\ -\ random(0,\ β\ ×\ log(compute\_time))}). Only one client recomputes at a time; others still serve stale value.
\item
  \textbf{Mutex / lock:} First cache-miss acquires a distributed lock (Redis \texttt{SET\ lock:key\ NX\ EX\ 5}). Others wait or serve stale.
\item
  \textbf{Background refresh:} Separate async job refreshes popular keys before they expire.
\item
  \textbf{Stale-while-revalidate:} Serve stale immediately; trigger async refresh; update cache when done.
\end{enumerate}

\textbf{Key decision: Stale-while-revalidate for most use cases} --- zero extra latency for users; background refresh handles expiry.

\subsubsection{6.7 Write Policies}\label{67-write-policies}

\setcodelang{}

\begin{verbatim}
Write-Through:   Client → Cache & DB (synchronous both)
                 + Consistent  − Higher write latency

Write-Back:      Client → Cache (immediate) → DB (async batch)
                 + Low latency  − Risk of data loss

Cache-Aside:     Client → DB directly; read-path populates cache
                 + DB authoritative  − Cold start cache misses
\end{verbatim}

\subsubsection{6.8 Node Failure Handling}\label{68-node-failure-handling}

\begin{itemize}
\tightlist
\item
  \textbf{Detection:} Gossip heartbeat; node marked \texttt{PFAIL} after missing 3 pings; cluster votes to mark \texttt{FAIL} (majority of masters agree).
\item
  \textbf{Failover:} Replica with smallest replication offset (least lag) wins election; becomes new master for the slot range within \textasciitilde1--3 s.
\item
  \textbf{Client impact:} During failover window, affected slots return \texttt{CLUSTERDOWN} errors \(\rightarrow\) client retries with backoff.
\item
  \textbf{Split brain prevention:} Requires majority of masters alive to accept writes. With 20 masters, lose 10 \(\rightarrow\) writes rejected (availability vs consistency trade-off).
\end{itemize}

\subsection{7. Bottlenecks \& Scaling}\label{7-bottlenecks--scaling-2}

\begin{longtable}[]{@{}
  >{\raggedright\arraybackslash}p{(\columnwidth - 2\tabcolsep) * \real{0.3133}}
  >{\raggedright\arraybackslash}p{(\columnwidth - 2\tabcolsep) * \real{0.6867}}@{}}
\toprule\noalign{}
\rowcolor{acc!16}
\begin{minipage}[b]{\linewidth}\raggedright
Bottleneck
\end{minipage} & \begin{minipage}[b]{\linewidth}\raggedright
Mitigation
\end{minipage} \\
\midrule\noalign{}
\endhead
\bottomrule\noalign{}
\endlastfoot
Hot key (single shard 100K QPS) & Key fan-out / local in-process cache \\
\rowcolor{zebra}
Memory exhaustion & Eviction policy; monitor \texttt{used\_memory\_rss}; add nodes \\
Network bandwidth (large values) & Compress values at client (snappy/lz4); max value size policy \\
\rowcolor{zebra}
Failover latency (1--3 s) & Keepalive + health checks; pre-elect backup masters \\
Replication lag & \texttt{WAIT\ 1\ 100} for critical writes; monitor \texttt{repl\_lag} metric \\
\rowcolor{zebra}
Cross-datacenter replication & Redis Cluster doesn\textquotesingle t span DCs natively; use CRDT-based solutions or geo-replicas with eventual consistency \\
\end{longtable}

\subsection{8. Trade-offs Summary}\label{8-trade-offs-summary-2}

\begin{longtable}[]{@{}
  >{\raggedright\arraybackslash}p{(\columnwidth - 6\tabcolsep) * \real{0.1959}}
  >{\raggedright\arraybackslash}p{(\columnwidth - 6\tabcolsep) * \real{0.1856}}
  >{\raggedright\arraybackslash}p{(\columnwidth - 6\tabcolsep) * \real{0.2371}}
  >{\raggedright\arraybackslash}p{(\columnwidth - 6\tabcolsep) * \real{0.3814}}@{}}
\toprule\noalign{}
\rowcolor{acc!16}
\begin{minipage}[b]{\linewidth}\raggedright
Decision
\end{minipage} & \begin{minipage}[b]{\linewidth}\raggedright
Option A
\end{minipage} & \begin{minipage}[b]{\linewidth}\raggedright
Option B
\end{minipage} & \begin{minipage}[b]{\linewidth}\raggedright
\textbf{Choice}
\end{minipage} \\
\midrule\noalign{}
\endhead
\bottomrule\noalign{}
\endlastfoot
Sharding strategy & Consistent hashing & Range-based sharding & \textbf{Consistent hashing} (smooth rebalance) \\
\rowcolor{zebra}
Eviction & LRU & LFU & \textbf{LFU} (scan-resistant) \\
Cache invalidation & TTL only & TTL + CDC events & \textbf{TTL + CDC} (real-time for critical) \\
\rowcolor{zebra}
Write policy & Write-through & Cache-aside & \textbf{Cache-aside} (DB authoritative) \\
Hot key & Replica reads & Client-side local cache & \textbf{Local cache} (zero network) \\
\rowcolor{zebra}
Stampede protection & Mutex lock & Stale-while-revalidate & \textbf{SWR} (no latency penalty) \\
Persistence & RDB only & AOF everysec & \textbf{AOF everysec} (\textless{} 1s loss) \\
\end{longtable}

\textbf{Interview takeaways:}

\begin{itemize}
\tightlist
\item
  \textbf{Consistent hashing with virtual nodes is essential} --- explain the 150-vnode trick for even distribution.
\item
  \textbf{Hot key and cache stampede are the two most common follow-up questions} --- have concrete solutions ready.
\item
  \textbf{Redis uses approximate LRU (sampling)} --- not true LRU. Know why.
\item
  Cache-aside is the most common pattern; write-through for critical financial data.
\end{itemize}

\section{4. Design a Web Crawler (Google-Scale)}\label{4-design-a-web-crawler-google-scale}

\subsection{1. Requirements}\label{1-requirements-3}

\textbf{Clarifying questions:}

\begin{itemize}
\tightlist
\item
  Full web crawl (all URLs), or focused crawl (specific domains/topics)?
\item
  How fresh does the data need to be? (daily recrawl vs. weekly vs. on-change)
\item
  Should we store raw HTML or parsed content?
\item
  Do we need to handle JavaScript-rendered pages (dynamic content)?
\item
  What\textquotesingle s the target crawl rate and storage budget?
\end{itemize}

\textbf{Functional Requirements:}

\begin{itemize}
\tightlist
\item
  Start from a set of seed URLs; discover new URLs via link extraction.
\item
  Download and store page content.
\item
  Respect \texttt{robots.txt} and crawl-delay directives.
\item
  Deduplicate URLs --- don\textquotesingle t crawl the same page twice (in a given window).
\item
  Freshness --- recrawl pages periodically based on update frequency.
\item
  Output: raw HTML (and/or extracted text) + metadata per URL.
\end{itemize}

\textbf{Non-Functional Requirements:}

\begin{itemize}
\tightlist
\item
  Scale: crawl 10 billion pages (Common Crawl scale).
\item
  Throughput: 1,000 pages/sec sustained (86M pages/day \(\rightarrow\) 10B pages in \textasciitilde116 days).
\item
  Distributed: hundreds of crawler nodes.
\item
  Politeness: max 1 req/sec per domain; honor robots.txt.
\item
  Durability: don\textquotesingle t re-crawl already-crawled URLs unless recrawl is due.
\item
  Freshness: high-value pages (news) re-crawled hourly; static pages monthly.
\end{itemize}

\subsection{2. Capacity Estimation}\label{2-capacity-estimation-3}

\setcodelang{}

\begin{verbatim}
Target pages           = 10 billion
Crawl rate             = 1,000 pages/sec
Time to crawl all      = 10B / 1,000 = 10,000,000 s ≈ 116 days  (1 full crawl cycle)
Peak rate needed       = 3,000–5,000 pages/sec (to finish faster / handle crawl failures)

Avg page size          = 100 KB HTML
Storage (HTML only)    = 10B × 100 KB = 1 PB
Compressed (~5:1 gzip) = ~200 TB

Links extracted per page = avg 50 outgoing links
Total URL queue entries  = 10B × 50 = 500B (dedup reduces this dramatically)
Seen URL set (dedup)     = 10B URLs × ~100 bytes/URL = 1 TB (in memory: impractical)
                           Bloom filter: 10B items, 1% FP → ~12 GB (fits in RAM!)

DNS lookups:
  1,000 pages/sec; avg 50 unique domains/sec
  DNS cache (TTL=300s) covers most; ~5 uncached DNS lookups/sec
\end{verbatim}

\subsection{3. API Design}\label{3-api-design-3}

\textbf{Internal service APIs (not user-facing):}

\setcodelang{Python}

\begin{Shaded}
\begin{Highlighting}[]
\CommentTok{\# URL Frontier Service}
\NormalTok{POST   }\OperatorTok{/}\NormalTok{frontier}\OperatorTok{/}\NormalTok{enqueue}
\NormalTok{       Body: \{ url, priority, source\_page, discovered\_at \}}
\NormalTok{       Returns: \{ queued: true }\OperatorTok{|}\NormalTok{ false (already seen) \}}

\NormalTok{GET    }\OperatorTok{/}\NormalTok{frontier}\OperatorTok{/}\BuiltInTok{next}\NormalTok{?crawler\_id}\OperatorTok{=\&}\NormalTok{batch\_size}\OperatorTok{=}\DecValTok{100}
\NormalTok{       Returns: [\{ url, fetch\_after\_ts \}]  (respects politeness delay)}

\NormalTok{POST   }\OperatorTok{/}\NormalTok{frontier}\OperatorTok{/}\NormalTok{done}
\NormalTok{       Body: \{ url, status\_code, crawled\_at, content\_hash \}}

\CommentTok{\# Content Store}
\NormalTok{PUT    }\OperatorTok{/}\NormalTok{content}\OperatorTok{/}\NormalTok{\{url\_hash\}}
\NormalTok{       Body: \{ raw\_html, headers, crawled\_at, content\_hash \}}

\NormalTok{GET    }\OperatorTok{/}\NormalTok{content}\OperatorTok{/}\NormalTok{\{url\_hash\}}

\CommentTok{\# Robots.txt Cache}
\NormalTok{GET    }\OperatorTok{/}\NormalTok{robots?domain}\OperatorTok{=}\NormalTok{example.com}
\NormalTok{       Returns: \{ allowed\_paths, disallowed\_paths, crawl\_delay\_s, sitemap\_urls \}}

\CommentTok{\# Recrawl Scheduler}
\NormalTok{POST   }\OperatorTok{/}\NormalTok{recrawl}\OperatorTok{/}\NormalTok{schedule}
\NormalTok{       Body: \{ url, next\_crawl\_at, priority \}}
\end{Highlighting}
\end{Shaded}

\subsection{4. Data Model}\label{4-data-model-3}

\textbf{URL Frontier (Priority Queue + Politeness Queue):}

\setcodelang{Python}

\begin{Shaded}
\begin{Highlighting}[]
\NormalTok{frontier\_urls (}
\NormalTok{  url\_hash        BIGINT  PRIMARY KEY,    }\OperatorTok{{-}{-}}\NormalTok{ xxHash64 of normalized URL}
\NormalTok{  url             TEXT,}
\NormalTok{  priority        FLOAT,                  }\OperatorTok{{-}{-}}\NormalTok{ higher }\OperatorTok{=}\NormalTok{ crawl sooner}
\NormalTok{  domain          TEXT,}
\NormalTok{  next\_fetch\_ts   TIMESTAMPTZ,            }\OperatorTok{{-}{-}}\NormalTok{ earliest time to fetch (politeness)}
\NormalTok{  status          ENUM(}\StringTok{\textquotesingle{}pending\textquotesingle{}}\NormalTok{,}\StringTok{\textquotesingle{}in\_flight\textquotesingle{}}\NormalTok{,}\StringTok{\textquotesingle{}done\textquotesingle{}}\NormalTok{,}\StringTok{\textquotesingle{}error\textquotesingle{}}\NormalTok{),}
\NormalTok{  retry\_count     INT DEFAULT }\DecValTok{0}
\NormalTok{)}
\NormalTok{Implemented }\ImportTok{as}\NormalTok{: Apache Kafka topics per priority band (high}\OperatorTok{/}\NormalTok{medium}\OperatorTok{/}\NormalTok{low)}
                \OperatorTok{+}\NormalTok{ Redis ZSETs per domain }\ControlFlowTok{for}\NormalTok{ politeness scheduling}
\end{Highlighting}
\end{Shaded}

\textbf{Crawled Pages (blob + metadata):}

\setcodelang{}

\begin{verbatim}
pages (
  url_hash        BIGINT PRIMARY KEY,
  url             TEXT,
  domain          TEXT,
  content_hash    BIGINT,        -- SimHash for near-duplicate detection
  crawled_at      TIMESTAMPTZ,
  http_status     INT,
  content_type    TEXT,
  raw_html_key    TEXT           -- S3 key for raw HTML blob
)
\end{verbatim}

\textbf{Seen URL Set (deduplication):}

\setcodelang{}

\begin{verbatim}
Bloom Filter (in-memory, distributed):
  - 10B URLs, 0.1% false positive rate → ~14.4 GB
  - Implementation: Redis BF.ADD / BF.EXISTS (RedisBloom module)
  - Fallback: persistent seen_urls table in Cassandra for exact check on BF hit
\end{verbatim}

\textbf{robots.txt Cache:}

\setcodelang{}

\begin{verbatim}
robots_cache (domain → RobotsTxt object, TTL=24h) in Redis
\end{verbatim}

\subsection{5. High-Level Architecture}\label{5-high-level-architecture-3}

\setcodelang{}

\begin{verbatim}
  Seed URLs
      │
      ▼
  ┌─────────────────────────────────────────────────────────────────────┐
  │                        URL Frontier Service                          │
  │                                                                      │
  │  Priority Queues (Kafka topics):                                     │
  │    high-priority:  news sites, .gov, .edu (crawl every 1 hr)       │
  │    medium:         major sites (crawl every 24 hr)                  │
  │    low:            long-tail / static (crawl every 30 days)         │
  │                                                                      │
  │  Politeness Scheduler (Redis ZSETs per domain):                     │
  │    ZADD domain:example.com <next_allowed_ts> <url>                  │
  │    Only release URLs when next_allowed_ts ≤ now                     │
  └──────────────────────────┬──────────────────────────────────────────┘
                             │  batch of URLs to fetch
                             ▼
  ┌─────────────────────────────────────────────────────────────────────┐
  │              Crawler Workers (100–500 nodes)                         │
  │                                                                      │
  │  Each worker:                                                        │
  │  1. Pull batch from Frontier (100 URLs)                             │
  │  2. DNS resolve (with local TTL cache)                              │
  │  3. Check robots.txt cache → skip disallowed URLs                   │
  │  4. HTTP GET (follow redirects, max 3)                              │
  │  5. Store raw HTML → S3                                             │
  │  6. Extract links → URL Extractor                                   │
  │  7. Report done → Frontier                                          │
  └──────────┬───────────────────────┬──────────────────────────────────┘
             │ raw HTML               │ extracted links
             ▼                        ▼
  ┌──────────────────┐      ┌─────────────────────────────────────────┐
  │   S3 Content      │      │          URL Extractor / Normalizer      │
  │   Store          │      │  - Resolve relative URLs                 │
  │   (raw HTML)     │      │  - Normalize (lowercase, strip fragment) │
  └──────────────────┘      │  - Dedup check (Bloom Filter in Redis)  │
                             │  - If new → enqueue in Frontier         │
                             └───────────────────┬─────────────────────┘
                                                 │
                                                 ▼
                             ┌───────────────────────────────────────┐
                             │   Bloom Filter (RedisBloom, 14 GB)    │
                             │   BF.EXISTS url_hash → seen? yes/no   │
                             └───────────────────────────────────────┘

  ┌─────────────────────────────────────────────────────────────────────┐
  │                     Downstream Consumers (via Kafka)                 │
  │  - Indexer: parse HTML → extract text → send to search index        │
  │  - Link Graph Builder: update PageRank graph                        │
  │  - Recrawl Scheduler: schedule next crawl based on change frequency │
  └─────────────────────────────────────────────────────────────────────┘
\end{verbatim}

\subsection{6. Deep Dives}\label{6-deep-dives-3}

\subsubsection{6.1 URL Frontier --- Priority + Politeness}\label{61-url-frontier--priority--politeness}

\textbf{Two-tier queue architecture (Mercator design):}

\setcodelang{}

\begin{verbatim}
Front-end queues (priority):      Back-end queues (politeness):
  Q_high  [news, trending]          Domain A: [url1, url3, url7]  → min 1s delay
  Q_med   [general web]             Domain B: [url2, url5]        → min 1s delay
  Q_low   [static, rarely updated]  Domain C: [url4, url6]        → min 5s delay

Selector: reads from front-end queues proportionally to priority.
          Routes each URL to its domain back-end queue.
          Back-end queue dispatcher waits for per-domain cooldown before releasing.
\end{verbatim}

\textbf{Priority assignment:}

\begin{itemize}
\tightlist
\item
  PageRank estimate of source page.
\item
  Domain authority (Alexa/Majestic rank).
\item
  Time since last crawl vs. estimated update frequency.
\item
  Explicit signals (sitemap priority, news feeds).
\end{itemize}

\textbf{Implementation:}

\begin{itemize}
\tightlist
\item
  Priority bands implemented as \textbf{Kafka topics} (\texttt{crawl.high}, \texttt{crawl.med}, \texttt{crawl.low}).
\item
  Politeness tracking: \textbf{Redis sorted set} per domain --- \texttt{ZADD\ domain:example.com\ \{next\_fetch\_ts\}\ \{url\}}. Poller does \texttt{ZRANGEBYSCORE\ domain:*\ -inf\ now\ LIMIT\ 10} every second.
\end{itemize}

\subsubsection{6.2 URL Deduplication --- Bloom Filter vs Exact Set}\label{62-url-deduplication--bloom-filter-vs-exact-set}

\begin{longtable}[]{@{}
  >{\raggedright\arraybackslash}p{(\columnwidth - 8\tabcolsep) * \real{0.2469}}
  >{\raggedright\arraybackslash}p{(\columnwidth - 8\tabcolsep) * \real{0.2469}}
  >{\raggedright\arraybackslash}p{(\columnwidth - 8\tabcolsep) * \real{0.1358}}
  >{\raggedright\arraybackslash}p{(\columnwidth - 8\tabcolsep) * \real{0.1852}}
  >{\raggedright\arraybackslash}p{(\columnwidth - 8\tabcolsep) * \real{0.1852}}@{}}
\toprule\noalign{}
\rowcolor{acc!16}
\begin{minipage}[b]{\linewidth}\raggedright
Approach
\end{minipage} & \begin{minipage}[b]{\linewidth}\raggedright
Storage for 10B URLs
\end{minipage} & \begin{minipage}[b]{\linewidth}\raggedright
Lookup time
\end{minipage} & \begin{minipage}[b]{\linewidth}\raggedright
False positives
\end{minipage} & \begin{minipage}[b]{\linewidth}\raggedright
False negatives
\end{minipage} \\
\midrule\noalign{}
\endhead
\bottomrule\noalign{}
\endlastfoot
\textbf{Bloom Filter (1\% FPR)} & \textasciitilde12 GB & O(k) hashes & 1\% (skip new URL) & 0\% \\
\rowcolor{zebra}
\textbf{Bloom Filter (0.1\% FPR)} & \textasciitilde14.4 GB & O(k) hashes & 0.1\% & 0\% \\
\textbf{Hash Set (in-memory)} & \textasciitilde1 TB (impractical) & O(1) & 0\% & 0\% \\
\rowcolor{zebra}
\textbf{Cassandra exact store} & \textasciitilde200 GB on disk & 1--5 ms & 0\% & 0\% \\
\textbf{Redis Set} & \textasciitilde600 GB & O(1) & 0\% & 0\% \\
\end{longtable}

\textbf{Key decision: Bloom Filter as first-pass dedup (RedisBloom, \textasciitilde14 GB), with Cassandra exact-set for conflict resolution on BF hit (false positive check).}

\begin{itemize}
\tightlist
\item
  BF says "seen" \(\rightarrow\) check Cassandra exact table (1\% of traffic \(\rightarrow\) Cassandra QPS manageable).
\item
  BF says "not seen" \(\rightarrow\) guaranteed new URL \(\rightarrow\) enqueue directly.
\end{itemize}

\textbf{URL normalization before hashing:}

\setcodelang{}

\begin{verbatim}
1. Lowercase scheme and host
2. Remove default port (:80, :443)
3. Sort query parameters alphabetically
4. Remove URL fragment (#anchor)
5. Resolve relative paths (/../ → /)
6. Strip tracking params (utm_source, fbclid, etc.)
7. Canonicalize: prefer canonical link rel tag in HTML
\end{verbatim}

\subsubsection{6.3 DNS Resolution}\label{63-dns-resolution}

\textbf{Problem:} 1,000 pages/sec from 1,000 different domains = 1,000 DNS lookups/sec. Public DNS resolvers (8.8.8.8) rate-limit at \textasciitilde few hundred QPS.

\textbf{Solutions:}

\begin{enumerate}
\def\labelenumi{\arabic{enumi}.}
\tightlist
\item
  \textbf{Local DNS cache per crawler node:} TTL-respecting LRU cache. Cache hit rate \textgreater{} 90\% for popular domains.
\item
  \textbf{Dedicated DNS resolver cluster:} Internal recursive DNS servers (unbound / BIND), scaled horizontally. Query rate: \textasciitilde50 uncached lookups/sec per node.
\item
  \textbf{Prefetch DNS:} When a new domain is discovered, pre-resolve asynchronously before it\textquotesingle s scheduled for crawl.
\item
  \textbf{DNS blacklist:} Skip private IPs (SSRF protection: 10.0.0.0/8, 192.168.0.0/16, 127.0.0.1).
\end{enumerate}

\subsubsection{6.4 robots.txt Compliance}\label{64-robotstxt-compliance}

\setcodelang{}

\begin{verbatim}
For domain example.com:
  1. Check Redis cache: GET robots:example.com
  2. Cache miss → fetch https://example.com/robots.txt
  3. Parse:
       User-agent: *
       Disallow: /private/
       Crawl-delay: 5
       Sitemap: https://example.com/sitemap.xml
  4. Store in Redis: SET robots:example.com <parsed_rules> EX 86400
  5. Before each URL fetch: check if path matches Disallow patterns

Politeness: max(1s_default, Crawl-delay_from_robots)
\end{verbatim}

\textbf{Gotchas:}

\begin{itemize}
\tightlist
\item
  robots.txt itself may not exist (HTTP 404 \(\rightarrow\) treat as "allow all").
\item
  robots.txt may be too large (truncate at 512 KB).
\item
  Different user-agent rules: use "Googlebot" or your crawler\textquotesingle s name.
\end{itemize}

\subsubsection{6.5 Trap Avoidance}\label{65-trap-avoidance}

\textbf{Spider traps} generate infinite URLs (e.g., \texttt{/calendar/next?date=...} \(\rightarrow\) infinite future dates).

Mitigation:

\begin{enumerate}
\def\labelenumi{\arabic{enumi}.}
\tightlist
\item
  \textbf{URL depth limit:} max 10 hops from seed URL.
\item
  \textbf{Per-domain URL count limit:} max 100K URLs per domain (override for news sites).
\item
  \textbf{URL path pattern detection:} if same path template seen \textgreater{} 1,000 times, blacklist pattern.
\item
  \textbf{Cycle detection:} if URL hash appears in path history (seen in URL chain) \(\rightarrow\) stop.
\item
  \textbf{Content hash dedup:} if page content SimHash is near-duplicate of already-crawled page \(\rightarrow\) skip storing, but do extract links (might have new ones).
\end{enumerate}

\subsubsection{6.6 Near-Duplicate Detection (SimHash)}\label{66-near-duplicate-detection-simhash}

\setcodelang{}

\begin{verbatim}
1. Extract tokens (words) from page HTML.
2. For each token, compute hash (64-bit).
3. For each bit position 0..63:
     if token_hash bit is 1 → add weight to column[bit]
     if token_hash bit is 0 → subtract weight
4. Final SimHash: bit[i] = 1 if column[i] > 0, else 0.

Two pages with SimHash Hamming distance ≤ 3 → near-duplicate.
Storage: 8 bytes per page × 10B pages = 80 GB → manageable in distributed store.
\end{verbatim}

\subsubsection{6.7 Distributed Coordination}\label{67-distributed-coordination}

\textbf{Problem:} 500 crawler nodes must not fetch the same URL concurrently.

\textbf{Solution:}

\begin{itemize}
\tightlist
\item
  URL Frontier is the single coordinator.
\item
  Each URL assigned to one crawler node (round-robin or URL-hash based routing).
\item
  Frontier marks URL \texttt{in\_flight} on dispatch; reverts to \texttt{pending} if heartbeat missing for 60 s (dead crawler).
\item
  No two nodes ever assigned the same URL from the frontier (frontier is authoritative).
\end{itemize}

\textbf{Crawler node health:}

\begin{itemize}
\tightlist
\item
  Periodic heartbeat to Frontier Service (every 10 s).
\item
  Dead node\textquotesingle s \texttt{in\_flight} URLs requeued automatically after 60 s timeout.
\end{itemize}

\subsubsection{6.8 Freshness \& Recrawl Scheduling}\label{68-freshness--recrawl-scheduling}

\begin{longtable}[]{@{}
  >{\raggedright\arraybackslash}p{(\columnwidth - 4\tabcolsep) * \real{0.2803}}
  >{\raggedright\arraybackslash}p{(\columnwidth - 4\tabcolsep) * \real{0.2102}}
  >{\raggedright\arraybackslash}p{(\columnwidth - 4\tabcolsep) * \real{0.5094}}@{}}
\toprule\noalign{}
\rowcolor{acc!16}
\begin{minipage}[b]{\linewidth}\raggedright
Page type
\end{minipage} & \begin{minipage}[b]{\linewidth}\raggedright
Crawl frequency
\end{minipage} & \begin{minipage}[b]{\linewidth}\raggedright
Signal
\end{minipage} \\
\midrule\noalign{}
\endhead
\bottomrule\noalign{}
\endlastfoot
News article (first 1h) & Every 15 min & Source is major news domain + recent publish \\
\rowcolor{zebra}
News article (1--24h) & Every 2 hr & Declining share velocity \\
Popular blog post & Daily & Sitemap \texttt{changefreq=daily} \\
\rowcolor{zebra}
Product page & Weekly & Sitemap \texttt{changefreq=weekly} \\
Static reference page & Monthly & Sitemap \texttt{changefreq=monthly} or inferred \\
\rowcolor{zebra}
Newly discovered URL & Once (ASAP) & Not yet crawled \\
\end{longtable}

\textbf{Adaptive recrawl:}

\begin{itemize}
\tightlist
\item
  Track \texttt{change\_rate\ =\ content\_hash\_changes\ /\ recrawls} per URL.
\item
  Increase frequency if page changes often; decrease if rarely changes.
\item
  \texttt{next\_crawl\_delay\ =\ max(min\_delay,\ last\_crawl\_interval\ ×\ (1\ -\ change\_rate\ +\ 0.1))}
\end{itemize}

\subsection{7. Bottlenecks \& Scaling}\label{7-bottlenecks--scaling-3}

\begin{longtable}[]{@{}
  >{\raggedright\arraybackslash}p{(\columnwidth - 2\tabcolsep) * \real{0.3554}}
  >{\raggedright\arraybackslash}p{(\columnwidth - 2\tabcolsep) * \real{0.6446}}@{}}
\toprule\noalign{}
\rowcolor{acc!16}
\begin{minipage}[b]{\linewidth}\raggedright
Bottleneck
\end{minipage} & \begin{minipage}[b]{\linewidth}\raggedright
Mitigation
\end{minipage} \\
\midrule\noalign{}
\endhead
\bottomrule\noalign{}
\endlastfoot
URL Frontier throughput & Partition Kafka topics by domain-hash; horizontal Frontier shards \\
\rowcolor{zebra}
Bloom filter accuracy drift & Use Counting Bloom Filter or periodic rebuild; monitor FP rate \\
DNS resolution rate & Dedicated internal DNS cluster + per-node LRU cache \\
\rowcolor{zebra}
S3 write throughput & Randomize S3 key prefix; use multipart upload for large HTML batches \\
robots.txt fetch overhead & Redis cache (TTL=24h); prefetch on domain discovery \\
\rowcolor{zebra}
Spider traps infinite loop & URL depth limit + per-domain URL count cap \\
Crawler node failure & Frontier requeues in\_flight URLs after heartbeat timeout \\
\rowcolor{zebra}
Hot domains (Reddit, Twitter) & Per-domain rate limit enforced strictly; queue backs up, not domain \\
\end{longtable}

\subsection{8. Trade-offs Summary}\label{8-trade-offs-summary-3}

\begin{longtable}[]{@{}
  >{\raggedright\arraybackslash}p{(\columnwidth - 6\tabcolsep) * \real{0.1926}}
  >{\raggedright\arraybackslash}p{(\columnwidth - 6\tabcolsep) * \real{0.2167}}
  >{\raggedright\arraybackslash}p{(\columnwidth - 6\tabcolsep) * \real{0.2408}}
  >{\raggedright\arraybackslash}p{(\columnwidth - 6\tabcolsep) * \real{0.3499}}@{}}
\toprule\noalign{}
\rowcolor{acc!16}
\begin{minipage}[b]{\linewidth}\raggedright
Decision
\end{minipage} & \begin{minipage}[b]{\linewidth}\raggedright
Option A
\end{minipage} & \begin{minipage}[b]{\linewidth}\raggedright
Option B
\end{minipage} & \begin{minipage}[b]{\linewidth}\raggedright
\textbf{Choice}
\end{minipage} \\
\midrule\noalign{}
\endhead
\bottomrule\noalign{}
\endlastfoot
URL dedup & Exact hash set (Cassandra) & Bloom filter + exact fallback & \textbf{BF + exact fallback} (14 GB vs 200 GB) \\
\rowcolor{zebra}
URL frontier storage & Database (Postgres) & Kafka topics + Redis ZSETs & \textbf{Kafka + Redis} (throughput) \\
DNS & Public resolvers (8.8.8.8) & Internal recursive DNS cluster & \textbf{Internal DNS} (rate limit avoidance) \\
\rowcolor{zebra}
Content storage & HDFS & S3 & \textbf{S3} (elastic, cheap, managed) \\
Near-duplicate detection & Full content MD5 comparison & SimHash (64-bit) & \textbf{SimHash} (O(1) compare, 8B/page) \\
\rowcolor{zebra}
JS rendering & Skip JS pages & Headless Chrome (Puppeteer) & \textbf{Skip first pass} (cost); headless for top-N sites \\
Crawl scheduling & Fixed interval & Adaptive (change-rate based) & \textbf{Adaptive} (better freshness/cost ratio) \\
\end{longtable}

\textbf{Interview takeaways:}

\begin{itemize}
\tightlist
\item
  \textbf{The URL Frontier is the heart of the crawler} --- the two-tier priority + politeness design (Mercator) is the canonical answer.
\item
  \textbf{Bloom filter for dedup at 10B scale} --- know the math: 10B items, 0.1\% FPR \(\rightarrow\) 14.4 GB. This is the key insight.
\item
  \textbf{Politeness is non-negotiable} --- not just ethical; aggressive crawling gets IP-banned, defeating the crawler.
\item
  \textbf{Distributed coordination:} URL Frontier as single source of truth prevents duplicate crawling across nodes.
\item
  \textbf{Spider trap avoidance:} URL depth limit + per-domain URL count cap are the two simple, effective defenses.
\end{itemize}

\emph{End of Part 2 --- Advanced System Design Problems}

\clearpage
\section{Visual Reference}
\noindent A set of figures distilling the key quantitative intuitions of this topic.\par\vspace{4pt}
\visualfigure{../figures/sdp-2/01-sharding-replication.pdf}{Sharding partitions data by key to scale writes across nodes; replication copies each shard for read scaling and failover. Most large datastores combine both.}
\end{document}
